\section{Conclusion}

Spaceflight does not act on the plant ionome in a way that generalises across
the two Veggie crops with archived elemental data. Lettuce shows no significant
change in any of ten elements; Tokyo Bekana shows \NSigCLR{} once dry-matter
dilution is accounted for, and a classifier trained on either crop fails on the
other. Apparent calcium and iron increases --- the specific accumulation
hypothesis this work set out to test --- are dilution artefacts and do not
survive compositional correction. Potassium, sodium and sulfur do.

Two contributions are methodological rather than biological, and may matter
more. A 1000-fold unit-labelling defect in OSD-655, present in the archived
data dictionary as well as the data files, would corrupt any pooled analysis
that trusted the stated units; it is caught by a plausibility gate that any
pipeline could run. And treating ICP-OES concentrations as compositional rather
than independent changes the conclusion for five of twelve elements.

On the question that motivated the work, dietary mineral overload from
spaceflight-grown produce is not achievable at plausible harvest scales ---
\OverloadFeLet\,kg of fresh tissue daily would be needed to reach the iron
limit. Sufficiency, not excess, is the operative question, and manganese is the
element to watch.

Progress now depends on access rather than analysis. Of \LSDAMeasured{} LSDA
experiments measuring a mineral or nutritional variable, \LsdaaccessOpen{} are
open and \LsdaaccessRequestrequired{} are request-gated. Two inexpensive
additions to future crop studies --- archived dry-matter fraction and substrate
elemental composition --- would convert several of this study's bounded
estimates into measurements.
